\documentclass[11pt]{article}
\usepackage[margin=1in]{geometry}
\usepackage{amsmath,amssymb,amsthm}
\usepackage[colorlinks=true,linkcolor=blue,citecolor=blue,urlcolor=blue]{hyperref}
\hypersetup{pdftitle={Exact zCDP for k-Ary Randomized Response},
  pdfauthor={Ji Ho Bae},pdfsubject={Exact zero-concentrated differential privacy of randomized response}}
\newtheorem{theorem}{Theorem}
\newtheorem{lemma}{Lemma}
\newtheorem{corollary}{Corollary}
\newcommand{\RR}{\operatorname{RR}}
\newcommand{\D}{\mathrm{D}}
\title{Exact zCDP for $k$-Ary Randomized Response}
\author{Ji Ho Bae\\LawAll AI\\
\href{mailto:jihobae@snu.ac.kr}{\texttt{jihobae@snu.ac.kr}}}
\date{}

\begin{document}
\maketitle

\begin{abstract}
We determine the least zero-concentrated differential privacy parameter of
$k$-ary randomized response for every $k\geq2$ and every privacy parameter
$\varepsilon\geq0$. The necessary condition for endpoint optimality obtained
by Harrison and Manurangsi is also sufficient, including the equality case.
When this condition fails, the normalized R\'enyi-divergence objective has a unique maximizer,
characterized by a scalar equation with an explicit finite bracket.
The proof expresses the objective as a secant slope of a function whose
curvature changes sign at most once. This characterization also yields
bisection with geometric error bounds for both the maximizing order and the
optimal zCDP parameter, in exact real arithmetic.
\end{abstract}

\section{Introduction and preliminaries}

Zero-concentrated differential privacy (zCDP), introduced by Bun and
Steinke~\cite{BS}, bounds R\'enyi divergence simultaneously at every order
$\alpha>1$. Determining the least admissible parameter therefore requires a
global optimization over the order. Harrison and Manurangsi
determine this parameter for several fundamental mechanisms. For $k$-ary
randomized response, they establish the endpoint formula for
$2\leq k\leq6$ and give a necessary condition for this formula to hold for general
$k$~\cite[Theorem~16 and Lemma~17]{HM}. They also note that their method
extends to moderate values of $k>6$ under more complicated
constraints~\cite[footnote~4]{HM}, while leaving a precise
characterization for large $k$ as an open problem~\cite[Section~7]{HM}.
We prove that their threshold is exact and characterize the optimum above it
by a unique root of an explicit scalar equation.

Fix an integer $k\geq2$ and write $[k]=\{1,\ldots,k\}$. For
$\varepsilon\geq0$, the mechanism $\RR_{k,\varepsilon}$ takes values in
$[k]$ and has output law $P_i$ on input $i$, where
\[
 P_i(j)=\Pr[\RR_{k,\varepsilon}(i)=j]
 =\begin{cases}
 e^\varepsilon/(e^\varepsilon+k-1),&j=i,\\
 1/(e^\varepsilon+k-1),&j\ne i.
 \end{cases}
\]
Every pair of distinct inputs is adjacent. For positive probability mass
functions $P,Q$ on $[k]$ and $\alpha>1$, the R\'enyi divergence is
\[
 \D_\alpha(P\Vert Q)
 =\frac{1}{\alpha-1}\log\sum_{j=1}^k
 P(j)^\alpha Q(j)^{1-\alpha}.
\]
The mechanism is $\rho$-zCDP, for $\rho\geq0$, if
\[
 \D_\alpha(P_i\Vert P_j)\leq\rho\alpha
 \qquad(i,j\in[k],\ \alpha>1).
\]
Let $\rho_*(k,\varepsilon)$ denote the least such $\rho$, and set
$q(\alpha)=\D_\alpha(P_1\Vert P_2)/\alpha$. All logarithms are natural.

For $Y\sim P_1$, define the privacy loss
$L=\log(P_1(Y)/P_2(Y))$ and its cumulant generating function
$\Lambda(s)=\log\mathbb{E}[e^{sL}]$. The variable $L$ has at most three
mass points. Atouani, Marchal, and Arbel~\cite{AMA} study optimal
sub-Gaussian variance proxies for three-mass distributions. For centered
privacy loss, this proxy is
$\sup_{s\ne0}2[\Lambda(s)-s\Lambda'(0)]/s^2$, whereas the least zCDP
parameter is $\sup_{s>0}\Lambda(s)/[s(s+1)]$. The latter leads to a
secant-slope problem with a positive base point.

\section{The exact parameter}

\begin{theorem}\label{thm:main}
For every integer $k\geq2$, $\rho_*(k,0)=0$. For $\varepsilon>0$, put
\[
 a=\frac{\varepsilon}{2},\qquad
 c=\frac{k-2}{2e^a},\qquad
 K(\varepsilon)=
 \frac{2(e^\varepsilon-1)(e^\varepsilon-1-\varepsilon)}
 {\varepsilon(e^\varepsilon+1)-2e^\varepsilon+2}.
\]
The denominator defining $K(\varepsilon)$ is strictly positive, and the
following alternatives hold.
\begin{enumerate}
\item If $k\leq K(\varepsilon)$, including equality, then
\begin{equation}\label{eq:endpoint}
 \rho_*(k,\varepsilon)=
 \frac{\varepsilon(e^\varepsilon-1)}{e^\varepsilon+k-1}.
\end{equation}
The function $q$ is strictly decreasing on $(1,\infty)$, and its supremum
is its limit as $\alpha\downarrow1$.
\item If $k>K(\varepsilon)$, the equation
\begin{equation}\label{eq:root}
 (x^2-a^2)\sinh x
 =2x(c+\cosh x)
 \log\frac{c+\cosh x}{c+\cosh a}
\end{equation}
has exactly one solution $x_*$ in the domain $x>a$. In this case
\begin{equation}\label{eq:interior}
 \rho_*(k,\varepsilon)=
 \frac{\varepsilon^2\sinh x_*}{2x_*(c+\cosh x_*)},
 \qquad
 \alpha_*=\frac12+\frac{x_*}{\varepsilon}
\end{equation}
and $\alpha_*>1$ is the unique maximizer of $q$. Moreover,
\begin{equation}\label{eq:bracket}
 a<x_*<U:=2\log\bigl(2(c+\cosh a)\bigr)
       =2\log(k-1+e^\varepsilon)-\varepsilon.
\end{equation}
\end{enumerate}
For every $\varepsilon\geq0$ and $\rho\geq0$, the mechanism
$\RR_{k,\varepsilon}$ is $\rho$-zCDP if and only if
$\rho\geq\rho_*(k,\varepsilon)$.
\end{theorem}

The threshold $K$ is due to Harrison and Manurangsi~\cite[Lemma~17]{HM}.
The endpoint $x=a$ also satisfies~\eqref{eq:root},
but lies outside the domain $x>a$.

\section{Proof of the characterization}

If $\varepsilon=0$, every output law is uniform, so all the divergences
vanish. Henceforth assume $\varepsilon>0$.

For any distinct inputs $i,j$, only the probabilities of symbols $i$ and $j$
differ. Substituting their values gives
\begin{equation}\label{eq:renyi}
 \D_\alpha(P_i\Vert P_j)
 =\frac{1}{\alpha-1}
 \log\frac{e^{\alpha\varepsilon}+e^{(1-\alpha)\varepsilon}+k-2}
 {e^\varepsilon+k-1}.
\end{equation}
This is the expression in~\cite[Proposition~20]{HM}. It is independent of
the ordered pair of distinct inputs; equal inputs have divergence zero.
Consequently
\begin{equation}\label{eq:sup}
 \rho_* =\sup_{\alpha>1}q(\alpha).
\end{equation}

Set
\[
 x=\varepsilon\left(\alpha-\frac12\right)>a,
 \qquad h(x)=\log(c+\cosh x),
 \qquad t=x^2,\quad t_0=a^2,
\]
and define
\[
 f(t)=h(\sqrt t)\quad(t>0),\qquad
 S(t)=\frac{f(t)-f(t_0)}{t-t_0}\quad(t>t_0).
\]
The numerator and denominator inside the logarithm
in~\eqref{eq:renyi} are $2e^a(c+\cosh x)$ and
$2e^a(c+\cosh a)$, respectively. Since
$\alpha(\alpha-1)=(x^2-a^2)/\varepsilon^2$, we obtain
\begin{equation}\label{eq:secant}
 q(\alpha)=\varepsilon^2 S(t),\qquad
 \rho_* =\varepsilon^2\sup_{t>t_0}S(t).
\end{equation}

The curvature pattern below is equivalent, under $u=\sinh x$, to the
monotonicity result of Guo, Luo, and Qi~\cite[Theorem~1]{GLQ}.

\begin{lemma}[Curvature]\label{lem:curvature}
For $c\geq0$, define
\[
 g(x)=x(1+c\cosh x)-\sinh x(c+\cosh x),\qquad x\geq0.
\]
For $x>0$,
\begin{equation}\label{eq:derivatives}
 f'(x^2)=\frac{\sinh x}{2x(c+\cosh x)},\qquad
 f''(x^2)=\frac{g(x)}{4x^3(c+\cosh x)^2}.
\end{equation}
If $c\leq2$, then $g(x)<0$ for all $x>0$. If $c>2$, there is a unique
$b>0$ such that $g(b)=0$, with $g>0$ on $(0,b)$ and $g<0$ on $(b,\infty)$.
\end{lemma}
\begin{proof}
Differentiation gives~\eqref{eq:derivatives} and
\begin{equation}\label{eq:gprime}
 g'(x)=\sinh x\bigl(cx-2\sinh x\bigr).
\end{equation}
The function $\sinh x/x$ increases strictly from $1$ to infinity on
$(0,\infty)$: its derivative has the sign of $x\cosh x-\sinh x$, which
vanishes at zero and has derivative $x\sinh x>0$.

For $c\leq2$, equation~\eqref{eq:gprime} gives $g'<0$ on $(0,\infty)$;
as $g(0)=0$, this proves the first case. For $c>2$, there is a unique
$r>0$ with $2\sinh r/r=c$. Thus $g$ increases strictly on $(0,r)$ and
decreases strictly on $(r,\infty)$, with $g(r)>0$. Finally,
\[
 \lim_{x\to\infty}e^{-2x}g(x)=-\frac14.
\]
It follows that $g$ has exactly one positive zero $b>r$, with the stated
signs.
\end{proof}

The interior-maximum argument below is an instance of Yang's
l'Hospital-type monotonicity rule~\cite[Theorem~8]{Yang}; the direct proof
also provides the explicit bracket $U$.

\begin{lemma}[Secant slopes]\label{lem:shape}
If $g(a)\leq0$, then $S$ is strictly decreasing on $(t_0,\infty)$.
If $g(a)>0$, then $S$ increases strictly up to a unique
$t_*>t_0$ and decreases strictly thereafter. In the latter case
$c>2$, $a<b<\sqrt{t_*}<U$, and $t_*$ is the unique zero in $(t_0,\infty)$ of
\begin{equation}\label{eq:J}
 J(t)=(t-t_0)f'(t)-f(t)+f(t_0).
\end{equation}
\end{lemma}
\begin{proof}
The function $J$ extends to $t_0$ with $J(t_0)=0$, and
\[
 S'(t)=\frac{J(t)}{(t-t_0)^2},\qquad
 J'(t)=(t-t_0)f''(t).
\]
If $g(a)\leq0$, Lemma~\ref{lem:curvature} implies $g(x)<0$ for every
$x>a$, including when $g(a)=0$. Hence $J'<0$ on $(t_0,\infty)$, so
$J<0$ and $S'<0$ there.

If $g(a)>0$, the same lemma gives $c>2$ and $a<b$. The function $J$
increases strictly on $(t_0,b^2)$ and decreases strictly on
$(b^2,\infty)$; in particular, $J(b^2)>0$. Put
$C=\log(2(c+\cosh a))$, so $U=2C$. Since
$2(c+\cosh a)>e^a$, we have $C>a>0$. For $x\geq a$, the inequalities
$h'(x)\leq1$ and $h(x)\geq x-\log2$ give
\[
 J(x^2)=\frac{x^2-a^2}{2x}h'(x)-h(x)+h(a)
 \leq-\frac{x}{2}-\frac{a^2}{2x}+C.
\]
Thus $J(U^2)\leq-a^2/(4C)<0$. The monotonicity of $J$ now gives exactly
one zero $t_*\in(b^2,U^2)$, with $J>0$ on $(t_0,t_*)$ and $J<0$ on
$(t_*,\infty)$. The formula for $S'$ proves the assertion.
\end{proof}

\begin{proof}[Proof of Theorem~\ref{thm:main}]
To identify the threshold, write
\[
 A=a\cosh a-\sinh a,\qquad B=\sinh a\cosh a-a.
\]
Both are strictly positive for $a>0$, since they vanish at zero and have
derivatives $a\sinh a$ and $2\sinh^2a$, respectively. As $g(a)=cA-B$,
\begin{align*}
 g(a)\leq0
 &\iff c\leq B/A\\
 &\iff k\leq2+2e^a\frac{\sinh a\cosh a-a}{a\cosh a-\sinh a}\\
 &\iff k\leq
 \frac{2(e^{2a}-1)(e^{2a}-1-2a)}
 {2a(e^{2a}+1)-2e^{2a}+2}=K(\varepsilon).
\end{align*}
The denominator is $4e^aA>0$. Thus this equivalence also includes equality.

If $k\leq K(\varepsilon)$, Lemma~\ref{lem:shape}
and~\eqref{eq:secant} give
\[
 \rho_* =\varepsilon^2 f'(a^2)
 =\frac{\varepsilon^2\sinh a}{2a(c+\cosh a)}
 =\frac{\varepsilon(e^\varepsilon-1)}{e^\varepsilon+k-1}.
\]
Both changes of variable are strictly increasing, so $q$ is strictly
decreasing and its supremum is approached as $\alpha\downarrow1$.

If $k>K(\varepsilon)$, let $x_*=\sqrt{t_*}$, where $t_*$ is supplied by
Lemma~\ref{lem:shape}. Multiplying $J(t_*)=0$ by
$2x_*(c+\cosh x_*)$ gives~\eqref{eq:root}; the lemma proves uniqueness
on $x>a$ and the bracket~\eqref{eq:bracket}. Moreover,
\[
 \sup_{t>t_0}S(t)=S(t_*)=f'(t_*),
\]
which proves~\eqref{eq:interior} and uniqueness of the maximizing order.

Finally, equation~\eqref{eq:sup} shows that every $\rho\geq\rho_*$
satisfies all the zCDP inequalities. If $\rho<\rho_*$, an inequality for
distinct inputs fails at $\alpha_*$ in the interior regime, or at an order
sufficiently close to $1$ in the endpoint regime. Together with the
$\varepsilon=0$ case, this proves the stated equivalence.
\end{proof}

\section{Computation}

Throughout this section, assume $\varepsilon>0$ and $k>K(\varepsilon)$.
Define, for
$x\geq a$,
\begin{equation}\label{eq:F}
 F(x)=\frac{x^2-a^2}{2x}\frac{\sinh x}{c+\cosh x}
       -\log\frac{c+\cosh x}{c+\cosh a}=J(x^2).
\end{equation}
Lemma~\ref{lem:shape} gives
\begin{equation}\label{eq:Fsign}
 F(x)>0\quad(a<x<x_*),\qquad
 F(x_*)=0,\qquad F(x)<0\quad(x>x_*).
\end{equation}
This sign pattern permits bisection from $(a,U)$, despite the additional
zero $F(a)=0$ at the excluded endpoint.

\begin{corollary}\label{cor:bisection}
Assume $\varepsilon>0$ and $k>K(\varepsilon)$. Set $L_0=a$, $R_0=U$,
and $m_n=(L_n+R_n)/2$. If $F(m_n)>0$, set
$(L_{n+1},R_{n+1})=(m_n,R_n)$; if $F(m_n)<0$, set
$(L_{n+1},R_{n+1})=(L_n,m_n)$; if $F(m_n)=0$, terminate with $x_*=m_n$.
Before exact termination,
\[
 L_n<x_*<R_n,\qquad R_n-L_n=(U-a)2^{-n}.
\]
The corresponding enclosure for $\alpha_*$ is
$\bigl(\tfrac12+L_n/\varepsilon,\tfrac12+R_n/\varepsilon\bigr)$.
The quantities
\begin{equation}\label{eq:valueinterval}
 \ell_n=\frac{\varepsilon^2 h'(L_n)}{2R_n},\qquad
 u_n=\frac{\varepsilon^2 h'(R_n)}{2L_n}
\end{equation}
satisfy $\ell_n\leq\rho_*\leq u_n$, converge to $\rho_*$, and obey
\begin{equation}\label{eq:error}
 u_n-\ell_n\leq(\varepsilon+2)(U-a)2^{-n}.
\end{equation}
Thus an enclosure of width at most $\delta>0$ is obtained within
\[
 \max\left\{0,\left\lceil
 \log_2\frac{(\varepsilon+2)(U-a)}{\delta}
 \right\rceil\right\}
\]
bisection steps, or the exact root is found earlier.
\end{corollary}
\begin{proof}
The bracket invariant follows from~\eqref{eq:Fsign}. Differentiation gives
\[
 h''(x)=\frac{1+c\cosh x}{(c+\cosh x)^2}>0.
\]
Thus $h'$ is increasing, and $L_n\leq x_*\leq R_n$
together with~\eqref{eq:interior} gives $\ell_n\leq\rho_*\leq u_n$.
The shrinking bracket and continuity give convergence.

For $x>0$, we have $0\leq h'(x)\leq1$ and $h''(x)\leq1$, the latter
following from
\[
 (c+\cosh x)^2-(1+c\cosh x)
 =c^2+c\cosh x+\sinh^2x\geq0.
\]
Writing $L=L_n$, $R=R_n$, and $w=R-L$, we therefore have
\begin{align*}
 u_n-\ell_n
 &=\frac{\varepsilon^2}{2}
 \left(\frac{h'(R)-h'(L)}{L}
       +h'(L)\left(\frac1L-\frac1R\right)\right)\\
 &\leq\frac{\varepsilon^2}{2}
 \left(\frac{w}{a}+\frac{w}{a^2}\right)
 =(\varepsilon+2)w.
\end{align*}
Substituting $w=(U-a)2^{-n}$ proves~\eqref{eq:error} and the iteration
bound.
\end{proof}

Corollary~\ref{cor:bisection} assumes exact evaluation of elementary real
functions and exact sign comparisons; its iteration bound is not
a bit-complexity bound. A rigorous numerical implementation
must also control rounding in the threshold comparison, the sign tests,
and the value enclosure.

\section*{AI use disclosure}

The author used AI tools (GPT-6 Astra Max and Opus 5.5 Max) for code review,
assistance with manuscript translation, and the development of counterexamples.
The author takes full responsibility for the code, the original manuscript,
the final translation, and all other related outputs.

\begin{thebibliography}{9}
\bibitem{AMA}
Soufiane Atouani, Olivier Marchal, and Julyan Arbel.
\newblock Optimal sub-Gaussian variance proxy for 3-mass distributions.
\newblock arXiv:2510.06132v1, 2025.
\newblock \url{https://arxiv.org/abs/2510.06132v1}.

\bibitem{BS}
Mark Bun and Thomas Steinke.
\newblock Concentrated Differential Privacy: Simplifications, Extensions,
and Lower Bounds.
\newblock In \emph{Theory of Cryptography (TCC 2016-B), Part I},
LNCS 9985, pages 635--658, 2016.
\newblock \href{https://doi.org/10.1007/978-3-662-53641-4_24}
{doi:10.1007/978-3-662-53641-4\_24}.
\newblock Full version: \url{https://arxiv.org/abs/1605.02065v1}.

\bibitem{GLQ}
Bai-Ni Guo, Qiu-Ming Luo, and Feng Qi.
\newblock Monotonicity results and inequalities for the inverse hyperbolic
sine function.
\newblock \emph{Journal of Inequalities and Applications}, 2013:536, 2013.
\newblock \href{https://doi.org/10.1186/1029-242X-2013-536}
{doi:10.1186/1029-242X-2013-536}.

\bibitem{HM}
Charlie Harrison and Pasin Manurangsi.
\newblock Exact zCDP Characterizations for Fundamental Differentially
Private Mechanisms.
\newblock In \emph{7th Symposium on Foundations of Responsible Computing
(FORC 2026)}, LIPIcs 368, Article 3, pages 3:1--3:18, 2026.
\newblock \href{https://doi.org/10.4230/LIPIcs.FORC.2026.3}
{doi:10.4230/LIPIcs.FORC.2026.3}.

\bibitem{Yang}
Zhen-Hang Yang.
\newblock A new way to prove L'Hospital Monotone Rules with applications.
\newblock arXiv:1409.6408v1, 2014.
\newblock \url{https://arxiv.org/abs/1409.6408v1}.
\end{thebibliography}
\end{document}
